\chapter{Traitement des boues des stations d'épuration urbaines}

L’objectif de ce module est d’acquérir des connaissances générales sur les boues de stations d’épuration, ainsi que sur les techniques d’incinération de ces boues dans une UIOM.



\section{D'où viennent les boues ?}


En France chaque personne utilise en moyenne 200 litres d’eau par jour pour divers usages domestiques. Ces eaux usées sont rejetées dans un réseau d’assainissement chargé de les transporter vers une station d’épuration. 

\begin{figure}[H]
\centering
\includegraphics[width=\textwidth]{Images/M11/M11_1.png}
\end{figure}

La fonction de la station d’épuration est de retirer les matières indésirables pour pouvoir
rejeter de l’eau propre dans le milieu naturel. Dans certains cas, l’eau épurée peut être 
réutilisée.
\vspace{0.3cm}

La majeure partie des matières retirées à l’eau se retrouve concentrée dans ce que l’on appelle les « boues urbaines ». Ces boues contiennent majoritairement de l’eau (plus de 70\%) et de la matière organique.
La présence de certaines substances métalliques, telles que le cadmium et le mercure (qui rentrent dans la composition des piles), pose des problèmes croissants pour l’épandage agricole. 
\vspace{0.3cm}

À partir de 2002, la mise en décharge de ces boues sera très probablement interdite. Le traitement des boues de stations d’épuration par incinération devient un sujet de plus en plus d’actualité. C’est pourquoi plusieurs de nos incinérateurs traitent des boues de stations d’épuration.



\textbf{Production de boues urbaines}
\vspace{0.3cm}


\begin{figure}[H]
\centering
\includegraphics[width=\textwidth]{Images/M11/M11_2.png}
\end{figure}


\textbf{Les boues industrielles}
\vspace{0.3cm}

Certaines industries produisent elles aussi des boues issues du traitement de leurs eaux usées. Suivant les cas, ces boues « industrielles » pourront être considérées comme des « déchets industriels banals » (par exemple : fabrication de papiers ou industrie agro-alimentaire).

\vspace{0.3cm}

Dans d’autre cas, la composition chimique de ces boues les classera parmi les « déchets industriels spéciaux » et elles devront être dirigées dans des unités de traitement adaptées (par exemple : traitements de surface, chimie).

 

\section{Les différentes filières d’élimination des boues urbaines}



\begin{figure}[H]
\centering
\includegraphics[width=0.5\textwidth]{Images/M11/M11_3.png}
\end{figure}


\subsection{La valorisation agricole}

\textbf{Épandage direct}
\vspace{0.3cm} 

La valorisation agricole des boues découle de pratiques anciennes. De plus, les boues ont en général une bonne valeur agronomique (présence de fertilisants, d’oligo-éléments et de chaux). Actuellement, l’épandage de boues concerne moins de 1\% des terrains agricoles exploités, contre plus de 90\% pour les fumiers.




\begin{figure}[H]
\centering
\includegraphics[width=\textwidth]{Images/M11/M11_4.png}
\end{figure}

\vspace{0.3cm}



\textbf{L’épandage après compostage}
\vspace{0.3cm}

Le compostage est une technique relativement ancienne qui est, entre autres, utilisée pour les ordures ménagères et les déchets verts. Pratiquement, le compostage correspond à une dégradation biologique des matières organiques. Le compostage de boues s’effectue en aérant un mélange de boues fraîches et de coproduits (agent structurant riche en fibres tel que les déchets verts ou les écorces de bois).
\vspace{0.3cm}

Par rapport à un épandage « classique », le compostage présente l’avantage de réduire la quantité de produit, de limiter les odeurs (stabilisation) et d’obtenir un produit sans risque pour la santé publique.
\vspace{0.3cm}

Il existe déjà des plates-formes de compostage de boues en France. Chaque année 40 000 tonnes d’OM (après tri sélectif de la fraction organique) et 10 000 tonnes de boues de stations d’épuration sont ainsi traitées et entrent dans un plan d’épandage. 


\subsection{Mise en décharge}

La réglementation exige une humidité inférieure à 70\% pour le stockage en CET de classe 2, à condition de ne pas présenter de nuisances pour l’environnement.
À partir de 2002, le stockage en décharge devra être limité aux déchets dits « ultimes », c’est-à-dire ne pouvant être valorisés « dans les conditions techniques et économiques du moment ». Dans ces circonstances, la plupart des boues urbaines ne pourront plus être admises en décharge. 


\section{Caractérisation des boues}

Les boues sont à définir selon trois critères : 

\begin{itemize}
\item une caractérisation physique qui donne des indications sur la structure des boues et leur comportement mécanique ;
\item  une caractérisation chimique (métaux, micro-polluants, teneur en eau, composition élémentaire) ;
\item  une caractérisation biologique qui permet d’évaluer les risques sanitaires.
\end{itemize}

La caractérisation des boues est très importante pour pouvoir définir correctement les filières d’élimination possibles. Par exemple, une boue dont les teneurs en micro-polluants organiques sont trop importantes devrait être dirigée soit vers l’incinération, soit vers la mise en décharge.


Plus précisément, on s’attachera à la liste des paramètres suivants :

\begin{itemize}
    \item La concentration en matières sèches (MS) exprimée en pourcentage du poids brut – siccité - (dans le cas des boues solides). Ce paramètre correspond à « tout ce qui n’est pas de l’eau dans une boue ». 
    Exemple : quand une boue a une siccité de 20\% (concentration en matières sèches de 20\%), cela signifie qu’elle a une teneur en eau de 80\%.

    \item La composition de la matière sèche, et notamment :
        \begin{itemize}
            \item la teneur en matières volatiles (MVS) exprimée en pourcentage de la MS ;
            \item la composition élémentaire en carbone et hydrogène qui entre en jeu dans le calcul du pouvoir calorifique inférieur (PCI) ;
            \item  la composition élémentaire en azote, phosphore, potassium, qui indique la valeur agronomique du produit ;
            \item  les teneurs en éléments traces métalliques (ETM), comme le cadmium, le plomb ou le mercure, déterminantes vis-à-vis de la possibilité d’épandre ou non les boues en agriculture.
        \end{itemize}
    \item La présence de micro-polluants organiques (due à la présence d'hydrocarbures et de détergents).
    \item La présence de micro-organismes pathogènes.
\end{itemize}


\begin{figure}[H]
\centering
\includegraphics[width=0.8\textwidth]{Images/M11/M11_5.png}
\end{figure}







\textbf{Les différents états physiques d'une boue}
\vspace{0.3cm}
En pratique, on utilise la classification suivante :


\begin{itemize}
\item Boues liquides : boues qui s'écoulent sous leur propre poids (siccité : 10\%).
\item Boues pâteuses : boues qui ne sont ni liquides ni solides (siccité : 10 à 30\%).
\item Boues solides : boues qui, entassées sur une hauteur de 1 mètre, forment une pente naturelle avec un angle de plus de 45° avec le sol.
\item Boues sèches : boues avec siccité > 60\%
\end{itemize}

 
 
 
 
\section{Les boues d'une station d'épuration}

\subsection{Traitement de l'eau}

Le schéma ci-dessous permet de visualiser plus précisément l’origine des boues dans une station d’épuration.

Le traiteur d’eaux pourra donc avoir jusqu’à trois grands types de boues à gérer s’il dispose d’un traitement primaire mécanique, d’un traitement biologique et d’un traitement physico-chimique (déphosphatation).


\vspace{0.3cm}
\textbf{Boues primaires "mécaniques", boues de décantation}
\vspace{0.3cm}

Elles sont le résultat d’une première décantation par gravité des eaux usées et sont donc composées en majorité de petits éléments minéraux.

\vspace{0.3cm}


\textbf{Boues biologiques}
\vspace{0.3cm}

En France, 98\% des stations urbaines utilisent des traitements par voie biologique.
Il existe plusieurs types de traitements biologiques des effluents : 

\begin{itemize} 
\item  simple circulation très lente dans des bassins directement creusés dans le sol 
(lagunage aéré) ; c’est la méthode la plus rustique ;
\item utilisation de bactéries pour décomposer la pollution organique contenue dans les eaux (on parle alors de « systèmes à boues activées ») ;
\item  épuration des eaux qui traversent une colonne composée de disques empilés (on parle de « disques biologiques »).
\end{itemize}

À chacune de ces solutions correspond un type de boue différent.

\begin{figure}[H]
\centering
\includegraphics[width=\textwidth]{Images/M11/M11_6.png}
\end{figure}


\vspace{0.3cm}
\textbf{Boues physico-chimiques}
\vspace{0.3 cm}

Les traitements physico-chimiques (ajouts de produits chimiques dans les eaux usées) concernent seulement 1,3\% des stations d’épuration urbaines françaises. Elles sont relativement difficiles à déshydrater.

\vspace{0.3cm}
\textbf{Bilan}
\vspace{0.3cm}

Le caractère commun de ces boues est de constituer un produit encore très liquide, de valeur énergétique généralement faible ou nulle. Certaines sont chimiquement inertes, mais celles qui proviennent de traitements biologiques peuvent fermenter et dégager des mauvaises odeurs. Afin de rendre possible leur manipulation, et éventuellement leur recyclage, le traiteur de boues doit définir une filière de traitement des boues adapté. 
Le traitement des boues a deux fonctions principales : 

\begin{itemize}
\item la réduction des teneurs en eau, c’est-à-dire augmenter leur concentration en matières sèches-MS ;
\item la stabilisation du produit, afin de limiter les nuisances olfactives et le nombre d’organismes pathogènes.
\end{itemize}


\subsection{Traitement des boues en vue de leur élimination}


\begin{figure}[H]
\centering
\includegraphics[width=0.6\textwidth]{Images/M11/M11_7.png}
\end{figure}





\textbf {Diminution de la teneur en eau}
\vspace{0.3cm}

L’augmentation de la siccité des boues commence par un épaississement qui a pour but de rendre possible une déshydratation mécanique ultérieure en effectuant une première réduction de la quantité d’eau contenue dans les boues.
Suite à l’épaississement, les boues vont devoir être stabilisées.


\vspace{0.3cm}
\textbf{Stabilisation}
\vspace{0.3cm}

L'objectif premier de la phase de stabilisation des boues de station d'épuration est de réduire les odeurs désagréables tout en limitant toute reprise de fermentation après traitement. En pratique, cette opération garantit que les boues sont plus respectueuses de l'environnement.

\vspace{0.3cm}
\textbf{Hygiénisation}
\vspace{0.3cm}

L’hygiénisation concerne principalement les boues destinées à l’épandage. Le but de cette opération est de réduire au maximum la toxicité des boues. 

\vspace{0.3cm}
\textbf{Déshydratation}
\vspace{0.3cm}

Une fois stabilisées, et selon l’exutoire défini préalablement, les boues vont subir une déshydratation mécanique qui va permettre de multiplier la siccité par trois ou quatre, voire dix pour certains filtres-presse. 
La siccité à obtenir après déshydratation doit être : 

\begin{itemize}
\item 15 - 30\% pour le compostage ;
\item 25 - 35\% pour assurer l'auto-combustibilité pour le traitement thermique dans un four spécifique ;
\item 15 - 30\% pour la co-incinération des boues pâteuses avec les ordures ménagères ;
\item plus de 30\% pour la mise en décharge.
\end{itemize}
 
En général, le coût de la déshydratation mécanique augmente avec le niveau d'efficacité requis. Sur cette base, les boues auront des taux de siccité très faibles pour les petites stations d'épuration rurales. 

\begin{figure}[H]
\centering
\includegraphics[width=0.8\textwidth]{Images/M11/M11_8.png}
\end{figure}

\begin{itemize}
\item La centrifugation directe : l’eau est séparée de la matière solide par une force centrifuge très importante. 
\item  Les filtres-presse : la boue est déshydratée par compression entre des plateaux.
\item  Les filtres à bande presseuse : la boue floculée (c’est-à-dire sous forme de flocons grâce à l’ajout de réactifs chimiques) est déshydratée par compression et cisaillement entre deux toiles.
\item  Le lit de séchage : les boues sont épandues sur un lit de graviers recouvert de sable. Par évaporation et infiltration, la boue s’épaissit. Cette technique est surtout utilisée sur des petites unités à la suite d’une stabilisation. 
\end{itemize}

Les principaux procédés utilisés pour une déshydratation mécanique des boues permettent d’atteindre des niveaux de siccité variables :

\begin{figure}[H]
\centering
\includegraphics[width=0.8\textwidth]{Images/M11/M11_9.png}
\end{figure}




 
\section{Incinération des boues}

À côté de la mise en décharge et de la valorisation agricole, l’incinération constitue la troisième voie principale d’élimination des boues de stations d’épuration.
Ce type de traitement peut être mis en place pour des boues de siccité supérieure à 15\%. Dans certains cas, un préséchage du produit intervient pour atteindre des siccités allant de 60 à 90\%. Ainsi, les boues peuvent être incinérées sur le site même de la station d’épuration ou en mélange avec des ordures ménagères. 
Dans les cas de l’incinération et de la co-incinération le bilan massique constaté est le suivant (pour des boues de siccité 25\% et de rapport MVS = 60\%) :



\begin{figure}[H]
\centering
\includegraphics[width =0.6\textwidth]{Images/M11/M11_10.png}
\end{figure}

\subsection{Incinération spécifique}

Dans la plupart des cas, l’incinération spécifique des boues de stations d’épuration urbaines est réalisée dans des fours à lit fluidisé, plus rarement dans des fours rotatifs et à soles étagées. Pour des raisons de prix de traitement des gaz de combustion notamment, ce type de filière est réservé aux stations d’épuration de grande taille.

La technique du lit fluidisé est la plus employée en France. Ceci s’explique par les avantages suivants :
\begin{itemize}
\item simplicité de la conception et de la réalisation ;
\item  facilité de mise en route, permettant facilement un fonctionnement discontinu ;
\item  faible débit de gaz de combustion ;
\item faible production d’oxyde d’azote (polluant présent dans les gaz de combustion) ;
\item  des gaz de combustion quasiment inodores en raison de leur température en sortie de four.
\end{itemize}


Le principe du four à lit fluidisé consiste à réaliser la combustion de la matière organique dans une enceinte verticale en acier garnie de réfractaire, de section cylindrique, dont le cœur est constitué d’un lit de sable fluidisé porté à haute température (environ 750 °C).
\vspace{0.3cm}

La fluidisation est obtenue en soufflant l’air de combustion sous le lit de sable à travers une grille équipée de tuyères en acier. La boue est introduite à l’intérieur du lit fluidisé en un ou plusieurs points selon la taille du four. Au-dessus du lit fluidisé, les gaz sont conduits dans une zone de post-réaction, « la zone revanche », dont le grand volume permet un temps de séjour très largement supérieur aux 2 secondes, à 850 °C, réglementaires.
Les points devant faire l’objet d’une attention particulière lors de la conception et de l’exploitation d’une telle installation sont :

\begin{itemize}
\item la récupération de chaleur dans les gaz de combustion ;
\item  l’épuration des gaz de combustion ;
\item  le devenir des cendres et des mâchefers.
\end{itemize}

\vspace{0.3cm}

Les deux procédés développés par Degremont et OTV sont respectivement NIRO®HTFB et 
PYROFLUID®.


\begin{figure}[H]
\centering
\includegraphics[width=0.8\textwidth]{Images/M11/M11_11.png}
\end{figure}





\vspace{0.3cm}
\textbf{Récupération de chaleur dans les gaz de combustion}
\vspace{0.3cm}

Dans la plupart des cas, la matière organique des boues possède un pouvoir calorifique insuffisant pour assurer à la fois l’évaporation de l’eau et l’élévation à 850 °C dans l’enceinte du four. Pour ne pas consommer de combustible d’appoint, il est donc nécessaire de récupérer une partie de l’énergie des gaz de combustion, soit par préchauffage de l’air de fluidisation dans un échangeur gaz de combustion-air, soit par préséchage des boues à des siccités supérieures à 40\% (ce qui revient à augmenter le PCI de la boue en réduisant la teneur en eau). La seconde méthode est plus efficace mais également plus coûteuse en terme d’investissement. 

\vspace{0.3cm}
\textbf{Épuration des gaz de combustion}
\vspace{0.3cm}

L’objectif est d’éliminer les polluants afin de respecter la réglementation relative aux émissions dans l’atmosphère. Les procédés de traitement des gaz de combustion sont variés :
\begin{itemize}
\item procédés secs : cyclones, filtres électrostatiques, filtres à manches, réacteurs de catalyse ;
\item procédés humides : laveurs à plateaux et à garnissage, filtres électrostatiques humides.
\end{itemize}


La conception de la chaîne et des équipements doit tenir compte d’une des spécificités de l’incinération sur four à lit fluidisé qui est l’envol de la totalité des matières minérales dans les gaz de combustion, conduisant à des teneurs en poussières bien supérieures à celles rencontrées pour les OM. Il faut également tenir compte de la destination finale des cendres et des résidus d’épuration des gaz de combustion.


\vspace{0.3cm}
\textbf{Devenir des cendres et des résidus d’épuration des gaz de combustion}
\vspace{0.3cm}

Depuis 1996, les cendres séparées par procédés uniquement physiques, sans addition de réactifs, peuvent, sous réserve de satisfaire à un certain nombre de tests, être valorisées en technique routière, dans la fabrication de ciments, être épandues en agriculture (cendres sèches) ou être évacuées en décharge de classe 2.
Par contre, les résidus de l’épuration physico-chimique (ajout de chaux), souvent présentés sous la forme d’un gâteau de filtre presse, doivent être évacués en CET de classe 1.
De ce fait, on cherche au maximum, à récupérer la matière minérale par des voies purement physiques de séparation, du type filtre électrostatique.
Les cendres volantes issues de l’incinération de boues en lit fluidisé semblent posséder des propriétés qui font défaut aux cendres d’OM : 

\begin{itemize}
\item les cendres de boues présentent des teneurs en éléments traces métalliques plus faibles que les cendres d’OM ;
\item  la fraction lixiviable des cendres volantes de boue est plus faible que celle des cendres d’OM.
\end{itemize}


\subsection{Co-incinération des boues avec les ordures ménagères}

\begin{itemize}
    \item  Co-incinération : incinération conjointe des ordures ménagères avec un autre déchet. Ce déchet peut bien entendu être des boues, mais aussi des DIB, des déchets d’activités de soins, etc.
    \item  Boues pâteuses : boues de siccité comprise entre 15 et 30 %.
    \item  Ratio brut boues/OM : rapport du tonnage de boues sur le tonnage d’OM durant une phase donnée de co-incinération.
\end{itemize}

Il faut dans ce cas prendre en compte les caractéristiques des boues, celles des ordures ménagères, les capacités de l’UIOM d’accueil et la réglementation.

\vspace{0.3cm}
\textbf{Co-incinération avec des boues pré-séchées}
\vspace{0.3cm}

Pour des stations d’épuration de taille importante, la co-incinération peut s’effectuer avec des boues préséchées à des siccités variant de 60 à 90\%. 

Il est à noter qu’à 65\% de siccité, les boues présentent un PCI voisin de celui des OM (soit environ 2 000 kcal/kg), ce qui facilite l’exploitation de l’unité. Les boues se présentent alors sous la forme de granulés, plus maniables et dont le stockage présente moins de contraintes. En effet, les techniques de séchage assurent également une très bonne stabilisation et hygiénisation du produit. Dans ces circonstances, le mélange se résume dans la plupart des cas à une introduction des boues en trémie par flux tendu. 
L’évolution technologique de la co-incinération de boues préséchées concerne plus les techniques de séchages préalables que le mode d’introduction proprement dit.

\vspace{0.3cm}
\textbf{Co-incinération avec des boues pâteuses}
\vspace{0.3cm}

Les expériences et les informations sur la co-incinération des boues pâteuses sont peu nombreuses. Quelques procédés sont validés et actuellement en marche sur site. C’est une technique en plein essor ayant déjà fait l’objet de plusieurs appels d’offres dans de grandes collectivités. À l’heure actuelle, des unités importantes sont en fonctionnement sur plusieurs UIOM.
Dans tous ces cas, la co-incinération concerne des boues pâteuses, ce qui évite un séchage préalable du produit et donc un coût d’exploitation trop élevé. 

\vspace{0.3cm}
\textbf{Séchage des boues}
\vspace{0.3cm}

L’objectif principal du séchage est de réduire fortement la masse (M) et le volume (V) des boues en provoquant l’évaporation de l’eau qu’elles contiennent :

Par ailleurs, cela permet d’hygiéniser le produit en détruisant les germes pathogènes.
Il existe trois modes de séchage :

\begin{figure}[H]
\centering
\includegraphics[width=0.6\textwidth]{Images/M11/M11_12.png}
\end{figure}


\vspace{0.3cm}
\textbf{Séchage direct}
\vspace{0.3 cm}

Le séchage, c’est-à-dire l’évaporation de l’eau contenue dans les boues, s’effectue grâce à une paroi chaude réchauffée au moyen d’un fluide (air chaud, par exemple) circulant dans une double enveloppe. Ce type de séchage concerne, entre autres, les unités de Rennes 
(Sobrec) et de Sète (Setom). 

\vspace{0.3cm}
\textbf{Séchage indirect}
\vspace{0.3cm}

Un gaz chaud et sec parcourt les boues et absorbe leur humidité. Ce type de séchage concerne par exemple les unités de Nice et de Quimper (Briec).



\begin{figure}[H]
\centering
\includegraphics[width=0.6\textwidth]{Images/M11/M11_13.png}
\end{figure}

\begin{figure}[H]
\centering
\includegraphics[width=0.4\textwidth]{Images/M11/M11_14.png}
\end{figure}

\vspace{0.3cm}
\textbf{Séchage mixte}
\vspace{0.3cm}

Mélange du séchage direct et indirect.

Le séchage des boues génère une buée qui correspond à l’eau contenue dans les boues qui s’est évaporée. Une partie de ces buées (partie liquide) est renvoyée dans une station d’épuration, l’autre (les incondensables) peut par exemple être recyclée dans le circuit d’air secondaire du four.


\begin{figure}[H]
\centering
\includegraphics[width=0.8\textwidth]{Images/M11/M11_15.png}
\end{figure}




\textbf{Principaux problèmes de fonctionnement d'un sécheur}
\vspace{0.3cm}

Les principaux types d’incidents rencontrés sont :
\begin{itemize}
\item les bourrages et les formations de dépôts dans les circuits ; 
\item  les départs d’incendies (dus à la présence de fines) en sortie de sécheurs ;
\item  les problèmes d’odeurs ;
\item une usure rapide du matériel.
\end{itemize}


\textbf{Modes d’introduction des boues pâteuses}
\vspace{0.3cm}

À ses débuts, la co-incinération se résumait à un stockage et à un simple mélange des boues avec les OM en fosse. La réglementation, appuyée par des contraintes d’exploitation croissantes, a poussé vers un stockage séparé et à des procédés où les boues ne rejoignent les ordures ménagères qu’à l’intérieur de l’enceinte du four.
Cela permet désormais une exploitation plus souple quant à la gestion des deux flux entrants. 


\begin{figure}[H]
\centering
\includegraphics[width=0.5\textwidth]{Images/M11/M11_16.png}
\end{figure}


\vspace{0.3cm}
\textbf{Introduction par émiettage mécanique}
\vspace{0.3cm}

Les boues sont émiettées grossièrement grâce à des coupelles rotatives et apposées sur la première zone de grilles, c’est-à-dire dans la zone de séchage. 
Grâce aux mouvements de la grille, cette technique permet de retenir les particules de boues au niveau de la zone de combustion, ce qui a pour effet de réduire l’impact de la co-incinération sur les Refiom (production moindre de cendres volantes). 
Ce mode d’introduction est utilisé par la technologie développée par 
Martin et Koch.

\vspace{0.3cm}

Cette méthode d'injection est utilisée dans les technologies développées par Martin et Koch.

\vspace{0.3cm}
\textbf{Introduction par extrusion}
\vspace{0.3cm}


Les boues sont filées dans les injecteurs et transformées en « boudins » d’une vingtaine de millimètres de diamètre. Ces boudins sèchent au cours de leur formation, puis se détachent sous l’effet de leur propre poids et tombent sur les grilles. De ce fait, ils sont mélangés aux ordures ménagères et incinérées. 
Le diamètre des agglomérats de boues étant important, le risque est de voir les boues cuites en surface et imbrûlées à l’intérieur. Ceci aurait pour effet d’augmenter considérablement la teneur en imbrûlés des mâchefers qui est par ailleurs limitée par la réglementation à moins de 5\%.
Cette technique est utilisée par les sociétés CNIM (en trémie avec les OM) et Degremont (procédé IC 850 installé sur la paroi du four).


\vspace{0.3cm}
\textbf{Introduction par pulvérisation}
\vspace{0.3cm}

À la sortie de la chambre de combustion
Cela consiste à pulvériser des boues pâteuses à contre-courant d’un flux de gaz ascendant à haute température, en l’occurrence les gaz de combustion provenant de l’incinération des OM. 
L’équipement utilisé se présente sous la forme d’une tour, installée au-dessus de la chambre de combustion du four. Les gaz de combustion chauds provenant de l’incinération des OM sont introduits à la base de cette tour, et un pulvérisateur situé à mi-hauteur projette de fines gouttelettes de boues vers le bas.
Au contact du flux de chaleur, la partie liquide de la boue se vaporise, et les matières sèches s’enflamment et brûlent instantanément. La majorité des boues est transformée en cendres volantes qui, si elles sont récupérées au niveau de l’électrofiltre, viennent augmenter la quantité de Refiom.
La société SGEE (procédé Ibisoc) développe cette technique.


\begin{figure}[H]
\centering
\includegraphics[width=0.6\textwidth]{Images/M11/M11_17.png}
\end{figure} 

Dans le four 
ce procédé constitue un compromis entre l’introduction par pulvérisation à la sortie de la chambre de combustion et l’introduction par extrusion. 
Les boues sont pulvérisées grâce à de l’air comprimé, ce qui permet de les transformer en particules très fines. Ces dernières tombent sur la grille (zone de séchage) et sont alors incinérées avec les ordures ménagères. Dans ces circonstances, la matière minérale contenue dans les boues se retrouve plutôt dans les mâchefers.
Ce mode d’injection est utilisé pour les procédés Inor (société Schwing) et Pyromix (Vivendi).


\begin{figure}[H]
\centering
\includegraphics[width=0.6\textwidth]{Images/M11/M11_18.png}
\end{figure}

\vspace{0.3cm}
\textbf{Contraintes de la co-incinération des boues avec les OM}
\vspace{0.3 cm}

Dans le cas d’un projet d’installation de co-incinération des boues avec les OM, différents aspects sont à prendre en compte.


\begin{figure}[H]
\centering
\includegraphics[width=0.6\textwidth]{Images/M11/M11_19.png}
\end{figure}

\vspace{0.3cm}
\textbf{Contexte socio-économique}
\vspace{0.3cm}

Même si la filière de valorisation en agriculture semble représenter la solution la moins chère, l’opinion publique n’est pas très favorable. Pour les communes qui ne peuvent pas épandre leurs boues, il faut trouver des solutions de remplacement.
Le choix de l’incinération spécifique est limité par un coût d’investissement très élevé. D’autre part, cette technique n’est pas encore satisfaisante en terme de valorisation énergétique. La co-incinération paraît donc une solution relativement bien adaptée, à condition qu’il existe une UIOM à proximité pouvant accueillir les quantités de boues et que les flux générés soient assez réguliers (en tonnage et sur la durée) pour rentabiliser l’investissement initial. Pour cela, on peut envisager des regroupements de stations dont les boues ont des caractéristiques à peu près homogènes (siccité, teneur en MVS).

\vspace{0.3cm}
\textbf{Contraintes techniques}
\vspace{0.3cm}

En plus des effets liés aux caractéristiques des produits entrants, l’impact de la co-incinération dépend des facteurs suivants :
\begin{itemize}
\item technologie du four ;
\item type de traitement des gaz de combustion ;
\item  conduite du four (respect des températures de combustion, volumes d’air primaire et secondaire…) ;
\item nature du procédé d’injection des boues.
\end{itemize}

\vspace{0.3cm}
Impact sur la charge de grille
\vspace{0.3cm}

Chaque fabricant précisera une charge maximale admissible (en kg/m2) en fonction de la technologie utilisée et de la capacité de l'incinérateur en termes de déchets ménagers . Une marge de sécurité sera généralement appliquée entre cette charge maximale admissible et la capacité optimale de l'incinérateur. Ajouter des boues aux ordures ménagères équivaut à augmenter la masse (par rapport à la seule incinération des ordures ménagères), qui doit être compatible avec la charge maximale sur la grille.



\vspace{0.3cm}
Impact sur les mâchefers
\vspace{0.3cm}

L’impact sur les mâchefers s’évalue par rapport aux limites fixées par la réglementation (circulaire de 1994).
La composition des boues étant différente de celle des OM, les mâchefers issus de la co-
incinération peuvent subir des modifications par rapport à une marche normale du four. Cela peut avoir pour conséquence de changer la classe du mâchefer (Valorisable, Maturable ou Stockable) et donc de modifier les coûts d’exploitation relatifs à l’évacuation de ces sous-produits (schématiquement 0, 100, 300 francs HT, respectivement pour V, M et S).
Le risque potentiel le plus élevé concerne la teneur en imbrûlés qui est limitée à 5\% par la réglementation.



\vspace{0.3cm}

La présence d'imbrûlés dans les mâchefers dépendra :
\vspace{0.3cm}

\begin{itemize}
\item du respect des températures de combustion ;
\item de la qualité de la répartition des boues dans le foyer afin d’éviter l’apparition de points froids localement (risque d’extinction de la flamme). Cela concerne plus particulièrement la co-incinération de boues pâteuses qui contiennent des teneurs en eau élevées ;
\item  de la taille moyenne des fragments de boue incinérés afin que le temps de séjour des boues dans le four soit compatible avec une bonne combustion du solide. De trop grosses particules risquent de provoquer « une cuisson à l’étouffée » qui se traduirait par la seule combustion de la périphérie.
Ceci signifie que le mode d’introduction des boues dans le four sera déterminant dans la qualité de la combustion. 
\end{itemize}

\vspace{0.3cm}
Impact sur le volume des gaz de combustion
\vspace{0.3cm}

Le débit des gaz de combustion est limité par le dimensionnement des installations de leur traitement et de la puissance du ventilateur de tirage qui maintient l’ensemble de l’installation en dépression. 
Cette dépression est nécessaire pour éviter les nuisances olfactives, les retours de flamme vers la fosse de stockage des OM, et les fuites de gaz de combustion vers l’atmosphère extérieure.


\vspace{0.3cm} 
La co-incinération de boues pâteuses peut entraîner une augmentation du volume des gaz de combustion produits. 
\vspace{0.3cm}

Cela s’explique par la teneur en eau relativement élevée qui génère de la vapeur lors de l’introduction dans le four.
La co-incinération avec des boues préséchées pose moins de problème de ce point de vue.


\vspace{0.3cm}
Impact sur REFIOM
\vspace{0.3cm}

Du fait de la teneur relativement élevée des boues en soufre (jusqu’à dix fois celle rencontrée pour les OM), la neutralisation des gaz de combustion, et donc du SO2, générera plus de Refiom. D’autre part, une augmentation de la production de cendres volantes, piégées au niveau de l’électrofiltre, se traduit également par une croissance de la production de Refiom.


\vspace{0.3cm}
Impact sur l'ensemble incinérateur-chaudière
\vspace{0.3cm} 

La co-incinération de boues avec les OM peut éventuellement conduire à une usure accrue de l’installation. Ces phénomènes peuvent se produire qu’après plusieurs mois, voire plusieurs années. La tenue générale du complexe four-chaudière lors d’une co-incinération dépend essentiellement des risques d’accrochages ou de projection de boues sur les parois. 

Ce type d’événement peut survenir :

\begin{itemize}
\item Dans le four : dans ce cas, ils peuvent être générés par une mauvaise maîtrise des températures de combustion (cendres volantes en fusion qui se solidifient au contact d’une paroi réfractaire plus froide) ou faire suite à une mauvaise conception du dispositif d’injection.

\item  Sur les éléments de la chaudière : comme précédemment, ils peuvent découler d’une mauvaise gestion des températures de combustion. Cela peut aboutir à une accélération des phénomènes de corrosion dont font l’objet les tubes de récupération de chaleur. Dans ces conditions, les performances d’échange thermique de la chaudière diminuent significativement et imposent à l’exploitant un ramonage plus fréquent.

\end{itemize}


 
\subsection{De la réception à l'incinération}

\textbf{Réception et stockage des boues}
\vspace{0.3cm}

La réception et le stockage des boues « brutes » se font en principe dans une fosse bétonnée indépendante de la fosse à OM. C’est le cas, par exemple, des usines de Rennes et de Sarcelles où les boues sont stockées avant d’être dirigées vers le sécheur à l’aide d’une pompe à boues. 

\vspace{0.3cm}
\textbf{Conditions de stockage des boues pâteuses}
\vspace{0.3cm}

Prévoir un système de dévoûtage
Dans tous les cas, l’élément de stockage principal doit être équipé d’un matériel assurant l’extraction des boues. En effet, ces produits ont tendance à « voûter », induisant ainsi des problèmes de vidange.

\vspace{0.3cm}
\textbf{Prévoir des ouvertures}
\vspace{0.3cm}

Bien que stabilisées, les boues continuent de fermenter. Cette fermentation génère la production de gaz méthane (CH4). Si aucune aération n’est prévue, le risque d’explosion est 
important. Il faut donc équiper les éléments de stockage d’« évents », c’est-à-dire d’ouvertures permettant au gaz de s’échapper dans l’atmosphère et évitant ainsi une surpression dangereuse.


\vspace{0.3cm}
\textbf{Réduire les nuisances olfactives}
\vspace{0.3cm}

La création des évents induit des dégagements d’odeurs nauséabondes dans l’atmosphère. Pour remédier à ces nuisances olfactives, on peut :
\begin{itemize}
\item raccorder les évents au système d’aspiration d’air primaire de l’UIOM ;
\item installer des filtres à charbon actif.
\end{itemize}

\vspace{0.3cm}
\textbf{Contraintes particulières de stockage des boues séchées}
\vspace{0.3cm}

La contrainte la plus importante en matière de stockage des boues séchées concerne les risques de départs d’incendie. 
Ainsi, de nombreux départs de feu ont été constatés au niveau de la trémie de stockage des fines. Ces risques d’inflammation de ces fines sont dans ce cas d’autant plus importants que les granulés de boues séchées sont à une température trop élevée en sortie de vis de refroidissement et que le taux de matières organiques dans les boues séchées est très élevé.

\vspace{0.3cm}
\textbf{Transport des boues vers l'incinérateur}
\vspace{0.3cm} 

Les boues peuvent être transférées de plusieurs manières :
\begin{itemize}
\item à l’aide du grappin lorsque les boues (en général séchées) sont stockées dans des fosses à ciel ouvert ;
\item  grâce à des convoyeurs à bandes depuis un silo de stockage ou un sécheur ;
\item  grâce à des pompes spécifiques aux matériaux épais, telles les pompes à béton (exemple : pompes à piston) dans le cas de boues pâteuses ;
\item grâce à un système pneumatique. Cette technique pose certains problèmes d’exploitation car on constate des bourrages et des abrasions au niveau des coudes ; 
\item  grâce à des élévateurs à godets.
\end{itemize}



\vspace{0.3cm}
\textbf{Exemple d'installation}
\vspace{0.3cm}


\begin{figure}[H]
\centering
\includegraphics[width=0.6\textwidth]{Images/M11/M11_20.png}
\end{figure}

\subsection{Fiche technique : calculs de PCI pour un mélange d'ordures ménagères + boues}


\vspace{0.3cm}
\textbf{PCI d'une boue}
\vspace{0.3cm}

Le PCI des boues, en kilocalories par kilogramme, peut être calculé comme suit :

\[
PCI_boue = 5 500 \times \text{(MVS content)} \times \text{(dryness)} - 600 \times \text{(1 - siccité)}
\]


\vspace{0.3cm}
\textbf{Exemple comment calculer le PCI des boues urbaines
(siccité 20\%, teneur en MVS 65\%)}
\vspace{0.3cm}

%\[
%PCI_boue = 	5.500 \times 0.65 \times 0.20 – 600 \times (1 – 0.20) = 235 \text{kcal/kg}
%\]

\begin{equation}
    PCI_{\text{boue}} = 5.500 \times 0.65 \times 0.20 - 600 \times (1 - 0.20) = 235 \, \text{kcal/kg}
\end{equation}


		
		Qui donne	qui donne	qui donne
		65 / 100	20 / 100	20 / 100



\vspace{0.3cm}
\textbf{PCI pour un mélange ordures ménagères + boues}
\vspace{0.3cm}

Commencer par calculer le pourcentage de boues et d'ordures ménagères dans le mélange incinéré :

\[
Rboue = \frac{\text{Tonnage boue}}{\text{(Tonnage ordures ménagères + tonnage boues)}}
\]

\[
RHshW = \frac{\text{Tonnage ordures ménagères}}{\text{((tonnage ordures ménagères + tonnage boues))}}
\]

Calculer ensuite le PCI du mélange en utilisant la formule suivante :

\[
NCVmixture = Rboue \times PCI_boue + RHshW \times PCIHshW
\]

\vspace{0.3cm}
\textbf{Exemple de calcul de NCV pour un mélange de boues + déchets ménagers}
\vspace{0.3cm}

\begin{itemize}

 \item  Compte tenu des boues ci-dessus (siccité 20\%, teneur en MVS 65\%), ce qui donne PCI = 235 kcal/kg
 \item  PCIHshW = 2 000 kcal/kg
 \item  1 tonne de boues brutes est incinérée pour 5 tonnes d'ordures ménagères
   
\end{itemize}


\[
Rboues = \frac{\text{1 t de boues}}{\text{((5 t d'ordures ménagères + 1 t de boues)}} = 0,17 = 17\%
\]

\[
RHshW = \frac{\text{ 5 t d'ordures ménagères}}{\text{((((5 t d'ordures ménagères + 1 t de boues)))}} = 0,83 = 83\%
\]

\[
\text{PCImélange} = 0,17 \times 235 + (0,83) \times 2000 = 1700 \text{kcal/kg}
\]

Ce qui donne	qui donne
17 / 100	83 / 100
\newpage
\subsection{LEXICON}

\vspace{0.3cm}
\textcolor{blue}{Bactéries}
\vspace{0.3cm}

Être vivant microscopique.

\vspace{0.3cm}
\textcolor{blue}{décantation}
\vspace{0.3cm}

Séparation naturelle des matières plus denses que l’eau qui se rassemblent dans le fond d’un récipient.

\vspace{0.3cm} 
\textcolor{blue}{Eaux usées} 
\vspace{0.3cm}

Eau rejetée après une utilisation donnée (exemple : eau de vaisselles).

\vspace{0,3cm} 
\textcolor{blue}{Éléments traces métalliques}
\vspace{0,3cm}

Composés que l’on appelle souvent à tort les « métaux lourds » (mercure, plomb, cadmium, etc.).


\vspace{0.3cm}
\textcolor{blue}{Micro-organismes organiques}
\vspace{0.3cm}


Composés toxiques dans les hydrocarbures, les détergents et les pesticides.

\vspace{0.3cm}
\textcolor{blue}{Épandage }
\vspace{0.3cm}

Opération qui consiste à répandre un produit à la surface d’un sol dans l’objectif d’une dégradation naturelle permettant un apport en matières fertilisantes.

\vspace{0.3cm}
\textcolor{blue}{Épuration des eaux}
\vspace{0.3cm}

Traitement des eaux usées avant rejet dans le milieu naturel (ou réutilisation).


\vspace{0.3cm}
\textcolor{blue}{Hygiénisation}
\vspace{0.3cm}

Opération dont le but est de détruire les micro-organismes pathogènes dans le but de préserver la santé publique.

\vspace{0.3cm}
\textcolor{blue}{Lagunage}
\vspace{0.3cm}

Technique d’épuration des eaux consistant à les faire circuler très lentement dans des bassins creusés directement dans le sol.

\vspace{0.3cm}
\textcolor{blue}{Micro-organismes}
\vspace{0.3cm}

Êtres vivants que l’on ne distingue pas à l’œil nu.

\vspace{0.3cm}
\textcolor{blue}{Micro-organismes pathogènes}
\vspace{0.3cm}

Petits êtres vivants qui sont à l’origine de nombreuses maladies.


\vspace{0.3cm}
\textcolor{blue}{Siccité}
\vspace{0.3cm}

Autre terme pour désigner la concentration en matières sèches contenues dans un échantillon.